\section{Fit to the Daisy data and genotype call (P.20, P.21)}
\exbox{Objective of the fit (P.20)}{The plant is compared with the Daisy runs by the objective $J(\theta)=\sum_k\bigl((m_k(\theta)-r_k)/\tau_k\bigr)^2$ over the relative yield, the transpiration share $\tau=T/(T+E_s)$, the mulch effect and five points of the E2 curve. The multipliers $\theta=(f_{\rm RUE},f_{kte},f_{qG})$ are searched one at a time with a step that is halved when no move helps. Using the assumed model values $m_k$, Daisy values $r_k$ and tolerances $\tau_k$ of Table~\ref{tab:J}, compute $J$.}\sol  Each term is $((m-r)/\tau)^2$. For the relative yield: $(0.86-0.90)/0.05=-0.8$, squared $0.64$. Summing the eight terms gives $J=4.447$. A one-parameter search tries $f_{\rm RUE}\to1.1f_{\rm RUE}$ and keeps the change if $J$ falls; if neither $\pm$ step helps the step is halved, within the bounds 0.6 to 1.6 (RUE), 0.5 to 2.2 ($k_{te}$) and 0.25 to 4 ($q_G$), in at most 45 model runs.\begin{table}[H]\centering\small\caption{Objective of the fit to the Daisy data.}\label{tab:J}
\begin{tabular}{lrrrr}\toprule
\hd{Target} & \hd{Model $m$} & \hd{Daisy $r$} & \hd{$\tau$} & \hd{$((m-r)/\tau)^2$}\\\midrule
relative yield & 0.86 & 0.90 & 0.05 & 0.640\\
transpiration share & 0.52 & 0.55 & 0.05 & 0.360\\
mulch effect on $E$ & 0.31 & 0.28 & 0.06 & 0.250\\
E2 point 1 & 0.97 & 1.00 & 0.04 & 0.563\\
E2 point 2 & 0.93 & 0.95 & 0.04 & 0.250\\
E2 point 3 & 0.88 & 0.90 & 0.04 & 0.250\\
E2 point 4 & 0.74 & 0.80 & 0.05 & 1.440\\
E2 point 5 & 0.55 & 0.50 & 0.06 & 0.694\\\bottomrule\end{tabular}\end{table}
\textbf{Reading.} The largest contributions are the points where the model departs by more than its tolerance; the sum is the quantity that the search minimises. A value of $J$ near the number of targets (8) means the model is off by about one tolerance on average.


\exbox{Genotype probabilities (P.21)}{The genotype call weights each genotype by $p_g\propto\exp\bigl(-(d_g-d_{\min})/(0.03+0.25d_{\min})\bigr)$, with $d_g$ the mean squared difference between its traits and the observed traits in units of the trait spread. Five genotypes have distances $0.40,0.55,0.20,0.90,1.60$ (assumed). Find the probabilities.}\sol  $d_{\min}=0.20$, so the temperature is $0.03+0.25(0.20)=0.08$. The weights are $e^{-(d-0.20)/0.08}$: for $d=0.40$, $e^{-2.5}=0.0821$; for $d=0.55$, $e^{-4.375}=0.0126$; $d=0.20$ gives 1; $d=0.90$ gives $e^{-8.75}=1.6\times10^{-4}$; $d=1.60$ gives $2.8\times10^{-7}$. The sum is 1.0948, so the probabilities are in Table~\ref{tab:geno}.\begin{table}[H]\centering\small\caption{Genotype call.}\label{tab:geno}
\begin{tabular}{lrrrr}\toprule
\hd{Genotype} & \hd{$d_g$} & \hd{$d_g-d_{\min}$} & \hd{Weight} & \hd{$p_g$}\\\midrule
G1 & 0.40 & 0.20 & 0.0821 & 0.075\\
G2 & 0.55 & 0.35 & 0.0126 & 0.011\\
G3 & 0.20 & 0.00 & 1.0000 & 0.913\\
G4 & 0.90 & 0.70 & 0.0002 & 0.000\\
G5 & 1.60 & 1.40 & 0.0000 & 0.000\\\bottomrule\end{tabular}\end{table}
The call is confident (about 91 per cent for G3) because the best distance is much smaller than the next; had the distances been $0.20$ and $0.22$ the weights would be $1$ and $0.78$ and the call would be uncertain. The temperature grows with $d_{\min}$, so a poor best fit never gives a sharp call.


\exbox{Solidity, DSI and STI (P.22)}{The projected shoot area of a plant is $A_{\rm proj}=420$ cm$^2$ and the convex hull around the leaves is $A_{\rm hull}=560$ cm$^2$. Four genotypes have grain under stress $Y_s$ and without stress $Y_p$ (Mg ha$^{-1}$, assumed) as in Table~\ref{tab:dsi}. Compute the solidity, the drought susceptibility index DSI $=(1-Y_s/Y_p)/(1-\bar Y_s/\bar Y_p)$ and the stress tolerance index STI $=Y_sY_p/\bar Y_p^2$.}\sol  Solidity $=A_{\rm proj}/A_{\rm hull}=0.750$. The means are $\bar Y_s=4.500$, $\bar Y_p=6.500$, so $1-\bar Y_s/\bar Y_p=0.3077$. For G2: $1-4.2/7.1=0.4085$, DSI $=0.4085/0.3077=1.327$; STI $=4.2\times7.1/6.500^2=0.706$.\begin{table}[H]\centering\small\caption{Drought susceptibility and stress tolerance indices.}\label{tab:dsi}
\begin{tabular}{lrrrr}\toprule
\hd{Genotype} & \hd{$Y_s$} & \hd{$Y_p$} & \hd{DSI} & \hd{STI}\\\midrule
G1 & 5.1 & 6.8 & 0.812 & 0.821\\
G2 & 4.2 & 7.1 & 1.327 & 0.706\\
G3 & 5.6 & 6.2 & 0.315 & 0.822\\
G4 & 3.1 & 5.9 & 1.542 & 0.433\\\bottomrule\end{tabular}\end{table}
\textbf{Reading.} DSI above 1 means the genotype loses more than the mean under stress (G2, G4); below 1 it is more tolerant (G3). STI rewards genotypes that are high in both environments: here G1 and G3 have the highest values.


\exbox{Depth of 95 per cent of the roots (P.23)}{Root length (cm) in the eight layers of 15 cm is $420,300,210,140,90,60,35,15$ (assumed). Find the depth $D_{95}$ at which the cumulative length reaches 95 per cent of the total, the root length density RLD $=L/(A\Delta z)$ of the top layer for a ground area $A=900$ cm$^2$, and the root-to-shoot ratio for $W_{\rm root}=18$ g and $W_{\rm shoot}=62$ g.}\sol  Total $=1270$ cm. Cumulative shares: 0.331, 0.567, 0.732, 0.843, 0.913, 0.961, 0.988, 1.000. The 95 per cent mark lies in layer 6 (cumulative 0.913 before it and 0.961 after); by linear interpolation inside the layer, $D_{95}=86.6$ cm. RLD of the top layer $=420/(900\times15)=0.0311$ cm cm$^{-3}$. Root-to-shoot ratio $=18/62=0.290$.


\section{The drought outlook and the stress calendar (P.24 to P.28)}
\exbox{Analogue years (P.25)}{Each of the 20 years of the weather has the rain of the days that have no forecast, $S_i=\sum_{k\in\rm future}P_i(k)$ (mm), $412,388,455,301,517,366,298,470,405,342,380,436,290,501,352,421,330,399,463,377$ (assumed). Find the dry, typical and wet analogue years $i_{\rm dry}=\arg\min S_i$, $i_{\rm wet}=\arg\max S_i$ and $i_{\rm typ}=$ median.}\sol  The minimum is $290$ (year 13), the maximum 517 (year 5). The sorted values are 290, 298, 301, 330, 342, 352, 366, 377, 380, 388, 399, 405, 412, 421, 436, 455, 463, 470, 501, 517; with 20 values the median is the mean of the 10th and 11th, $(388+399)/2=393.5$ mm, and the typical year is the one nearest to it: year 2 ($S=388$). The three futures are real years of the record, so each has its own dry spells in the order in which they happened; they are not a forecast.


\exbox{Periods of drought stress (P.26, P.27)}{The root-zone water is $a=\sum f_i\theta_i/\sum f_i$, the FAO-56 stress coefficient is $K_s=\min(1,a/(1-p))$ and a day is a stress day if $a<1-p$ ($p=0.55$ for maize and wheat). The root-zone values of 16 days are in Table~\ref{tab:ks}. Find the stress days, join gaps of at most 2 days, and report the periods (kept if $n_j\ge2$) with $a_{\min}$ and mean $K_s$.}\sol  The threshold is $1-p=0.45$. Stress days: those with $a<0.45$. The raw runs are [4,6], [8,10], [15,16]. Gaps of at most 2 days are filled, which merges these into [4,10], [15,16]. Periods kept: days 4 to 10 ($n=7$, $a_{\min}=0.33$, mean $K_s=0.867$); days 15 to 16 ($n=2$, $a_{\min}=0.35$, mean $K_s=0.822$).\begin{table}[H]\centering\small\caption{Root-zone water and the stress coefficient.}\label{tab:ks}
\begin{tabular}{lrrr}\toprule
\hd{Day} & \hd{$a$} & \hd{$K_s$} & \hd{Stress}\\\midrule
1 & 0.62 & 1.000 & -\\
2 & 0.55 & 1.000 & -\\
3 & 0.48 & 1.000 & -\\
4 & 0.40 & 0.889 & stress\\
5 & 0.38 & 0.844 & stress\\
6 & 0.44 & 0.978 & stress\\
7 & 0.52 & 1.000 & -\\
8 & 0.36 & 0.800 & stress\\
9 & 0.33 & 0.733 & stress\\
10 & 0.37 & 0.822 & stress\\
11 & 0.50 & 1.000 & -\\
12 & 0.60 & 1.000 & -\\
13 & 0.58 & 1.000 & -\\
14 & 0.46 & 1.000 & -\\
15 & 0.39 & 0.867 & stress\\
16 & 0.35 & 0.778 & stress\\\bottomrule\end{tabular}\end{table}


\exbox{Risk class and what the roots suffer (P.26, P.28)}{The risk class is set by $R=Y_{\rm rainfed}/Y_{\rm irrigated}$: low if $R\ge0.90$, moderate if $0.75\le R<0.90$, high if $0.50\le R<0.75$, severe below 0.50. For a dry future the rainfed grain is $3.1$ Mg ha$^{-1}$ and the irrigated $6.4$; typical $5.6/6.4$; wet $6.2/6.4$. Classify them, and compute the root measures of one period (P.28): root growth $\Delta L_k=(0.8,0.7,0.5,0.4)$ cm d$^{-1}\times$ plant against the bench plant $(1.0,1.0,1.0,0.9)$, deepest root $D=62$ cm against bench $70$ cm, and root lengths $L_i=(120,90,70,40,20)$ cm in layers with $\theta_i=(0.7,0.55,0.45,0.30,0.25)$.}\sol  $R=3.1/6.4=0.484$: severe; $5.6/6.4=0.875$: moderate; $6.2/6.4=0.969$: low. P.28: $G=100(0.8+0.7+0.5+0.4)/(1.0+1.0+1.0+0.9)=100(2.4)/(3.9)=61.5$ per cent (the roots grew about 62 per cent as much as the watered plant); $\Delta D=70-62=8$ cm shallower; layers with $\theta<0.5$ are layers 3, 4, 5: $F=100(70+40+20)/(120+90+70+40+20)=100(130)/340=38.2$ per cent of the root length lies in layers holding less than half of their plant-available water.


\section{Volumetric phenotyping (P2): rebuilding the scan}
Layer P2 measures the plant of P1 as a three-dimensional scanner would: views from a turning camera, calibration spheres, registration, a voxel cloud and traits.

\exbox{Joining two views by the spheres}{A calibration ring of six spheres of centres at radius $0.35$ m and angles $0,60,\dots,300$\textdegree\ is seen by view 2 turned by $60$\textdegree\ and shifted by $(0.020,-0.010)$ m, with centre-fit noise of $1$ mm. Four spheres (those at $0,60,120,180$\textdegree) are seen by both views. Estimate the rigid transform from view 2 back to view 1 (Procrustes/Kabsch) and its error.}\sol  1. Centroids: of the four observed centres $c_A=(-0.1109,0.0658)$ m, of the reference centres $c_B=(0.0000,0.1516)$ m. 2. Cross-covariance $H=\sum(A_i-c_A)(B_i-c_B)^T$; its singular value decomposition $H=U\Sigma V^T$ gives the rotation $R=V\,\mathrm{diag}(1,\det(VU^T))U^T$, with $\det(VU^T)=+1$ here so that $R$ is a proper rotation. 3. The rotation angle is $-60.07$\textdegree\ (truth: $-60$\textdegree\ to undo the turn), an error of $0.067$\textdegree. 4. Translation $t=c_B-Rc_A=(-0.0017,0.0226)$ m. 5. The residual after alignment has RMS $0.89$ mm, about the centre-fit noise of 1 mm. At the crown radius of $0.4$ m the angle error moves a point by $0.4\times0.00116=0.47$ mm.


\exbox{Accumulated registration error}{The transforms are chained to view 1: $R_0,\ R_0R_1,\ R_0R_1R_2,\dots$, so that the error of each link adds. If every link has a turn error of $\delta=0.3$\textdegree\ (a calibrated rig), compute the accumulated angle and the displacement at a point $0.4$ m from the axis for views $1$ to $6$; compare with the case of independent random errors that add in quadrature.}\sol  If the errors are all of the same sign the angle after $k$ links is $k\delta$ and the displacement is $r\,k\delta$ (in radians). For $k=6$: $1.8$\textdegree\ $=0.0314$ rad, displacement $0.4\times0.0314=12.6$ mm. If the link errors are independent with standard deviation $\delta$, the accumulated standard deviation is $\sqrt k\,\delta$, $0.73$\textdegree, so $5.1$ mm at $k=6$ (Table~\ref{tab:chain}). The last views drift most; the ICP refinement on the joined cloud is there to reduce this.\begin{table}[H]\centering\small\caption{Accumulated registration error at 0.4 m.}\label{tab:chain}
\begin{tabular}{lrrr}\toprule
\hd{View $k$} & \hd{Worst case (deg)} & \hd{Worst case (mm)} & \hd{Random (mm)}\\\midrule
1 & 0.30 & 2.09 & 2.09\\
2 & 0.60 & 4.19 & 2.96\\
3 & 0.90 & 6.28 & 3.63\\
4 & 1.20 & 8.38 & 4.19\\
5 & 1.50 & 10.47 & 4.68\\
6 & 1.80 & 12.57 & 5.13\\\bottomrule\end{tabular}\end{table}


\exbox{Height, crown, volume and leaf area from a voxel cloud}{A plant is modelled as a prolate ellipsoid with semi-axes $a=b=0.15$ m and $c=0.30$ m, standing on the ground (centre at $z=0.30$ m). A scan of 60,000 points is placed in voxels of $v=0.02$ m. The application measures: height = 99.8th percentile of $z$; crown width = longest horizontal span of the convex hull; canopy volume = occupied cubes; leaf area = surface voxels $\times v^2/1.5$. Compute them and compare with the truth.}\sol  Height: the 99.8th percentile of the point heights is 0.599 m (truth $0.600$ m; the percentile discards the top 0.2 per cent of points, so it is slightly below the apex). Crown width: the hull of the points projected on the ground has longest span 0.300 m (truth $2a=0.300$ m). Volume: the true volume is $\tfrac43\pi abc=28.27$ litres (0.02827 m$^3$); counting the cubes of $0.02$ m whose centres lie inside gives 3528 cubes, $V=28.22$ litres, an error of -0.2 per cent. Leaf area: the cloud occupies 1648 surface voxels, so $LA=1648\times0.02^2/1.5=0.439$ m$^2$ against the true surface of the ellipsoid 0.483 m$^2$ (Knud Thomsen approximation), a ratio of 0.91. The factor 1.5 is the application's convention; the voxel count depends on the voxel size (a coarse voxel counts the thickness of the leaf), as the manual states. The leaf area index for a ground area of $0.25$ m$^2$ is $LAI=LA/A=1.76$.


\exbox{Leaf area against the voxel size}{Repeat the leaf-area measurement for voxel sizes 0.01, 0.02, 0.04 and 0.08 m with the same cloud.}\sol  The number of surface voxels falls roughly as $v^{-2}$ so that $N v^2$ would be constant for a thin surface; a coarser voxel also counts the thickness, so the estimate departs from the true area. Table~\ref{tab:vox} shows the area and its ratio to the true surface (0.483 m$^2$).\begin{table}[H]\centering\small\caption{Leaf area against voxel size.}\label{tab:vox}
\begin{tabular}{lrrr}\toprule
\hd{Voxel (m)} & \hd{Surface voxels} & \hd{$LA$ (m$^2$)} & \hd{$LA$/true}\\\midrule
0.01 & 6194 & 0.413 & 0.85\\
0.02 & 1648 & 0.439 & 0.91\\
0.04 & 420 & 0.448 & 0.93\\
0.08 & 96 & 0.410 & 0.85\\\bottomrule\end{tabular}\end{table}
\textbf{Reading.} The estimate is not corrected for the voxel size in the application, and the error shown there is honest about it: a scanner's resolution is part of the measurement.


\exbox{Volumetric stress index}{The volumetric stress index is $\mathrm{VSI}=V_{\rm drought}/V_{\rm project}$. For a drought plant with canopy volume $18.4$ litres and a project plant with $26.9$ litres (assumed), and for scans every 10 days with volumes (drought, project) of $(2.1,2.3),(5.8,7.2),(11.9,15.8),(18.4,26.9)$ litres, compute the VSI each time and the transpiration per litre of canopy if the drought plant transpired 9.6 L and the project plant 14.2 L by the last date.}\sol  $\mathrm{VSI}=2.1/2.3=0.913$; $5.8/7.2=0.806$; $11.9/15.8=0.753$; $18.4/26.9=0.684$. The drought plant ends 32 per cent smaller. Transpiration per litre of canopy: drought $9.6/18.4=0.522$ L per litre; project $14.2/26.9=0.528$ L per litre; the two are within 1 per cent, so the drought plant is smaller in proportion to the water it used: it did not lose efficiency, it lost size.

